% p.14 — الباب السابع في ابواب الثلاثي المزيد والرباعي المجرد والمزيد
% Layout: the right-hand column carries group labels written vertically, each spanning the rows
% of its group (\vlabel); the one-row groups 21 and 22 have their labels set horizontally.
% The first block (abwāb 1–6 of the mujarrad) is a separate 3-row grid without column headers:
% [label] | # | ماضي | مضارع | مثال || # | ماضي | مضارع | مثال.
\mspage{14}
\nohairlines
\bab{الباب السابع في ابواب الثلاثي المزيد والرباعي المجرد والمزيد}
\begin{sarf}{|c|c|c|c|c|c|c|c|c|}
\vgrey & 1 & فَعَلَ & يَفْعَلُ & فَتَحَ & 4 & فَعِلَ & يَفْعَلُ & سَمِعَ\\
\vgrey & 2 & فَعَلَ & يَفْعِلُ & ضَرَبَ & 5 & فَعِلَ & يَفْعِلُ & حَسِبَ\\
\vlabel{3}{ثُلَاثِىْ}{مُجَرَّدٌ} & 3 & فَعَلَ & يَفْعُلُ & نَصَرَ & 6 & فَعُلَ & يَفْعُلُ & كَرُمَ\\
\end{sarf}
\begin{sarf}{|c|c|c|c|c|c|c|c|}
\thead{} & \vgrey & \thead{وزن اصلى} & \thead{وزن زائد} & \thead{مثال اصلى} & \thead{مثال زائد} & \thead{اصل} & \thead{\unclear{مزيد}}\\\hline
\vgrey & 7 & فَعَلَ & أَفْعَلَ & حَمَلَ & أَحْمَلَ & حَسَنَ & أَحْسَنَ\\
\vgrey & 8 & فَعَلَ & فَعَّلَ & حَمَلَ & حَمَّلَ & بَلَغَ & بَلَّغَ\\
\vlabel{3}{ثُلَاثِىْ مَزِيْدٌ}{بِحَرْفٍ} & 9 & فَعَلَ & فَاعَلَ & حَمَلَ & حَامَلَ & قَتَلَ & قَاتَلَ\\\hline
\vgrey & 10 & فَعَلَ & تَفَعَّلَ & حَمَلَ & تَحَمَّلَ & عَلِمَ & تَعَلَّمَ\\
\vgrey & 11 & فَعَلَ & تَفَاعَلَ & حَمَلَ & تَحَامَلَ & ثَقُلَ & تَثَاقَلَ\\
\vgrey & 12 & فَعَلَ & اِفْتَعَلَ & حَمَلَ & اِحْتَمَلَ & جَمَعَ & اِجْتَمَعَ\\
\vgrey & 13 & فَعَلَ & اِنْفَعَلَ & حَمَلَ & اِنْحَمَلَ & كَسَرَ & اِنْكَسَرَ\\
\vlabel{5}{ثُلَاثِىْ مَزِيْدٌ}{بِحَرْفَيْنِ} & 14 & فَعَلَ & اِفْعَلَّ & حَمَلَ & اِحْمَلَّ & حَمَرَ & اِحْمَرَّ\\\hline
\vgrey & 15 & فَعَلَ & اِسْتَفْعَلَ & حَمَلَ & اِسْتَحْمَلَ & كَرُمَ & اِسْتَكْرَمَ\\
\vgrey & 16 & فَعَلَ & اِفَّعَّلَ & حَمَلَ & اِحَّمَّلَ & زَمَلَ & اِزَّمَّلَ\\
\vgrey & 17 & فَعَلَ & اِفَّاعَلَ & حَمَلَ & اِحَّامَلَ & دَرَكَ & اِدَّارَكَ\\
\vgrey & 18 & فَعَلَ & اِفْعَالَّ & حَمَلَ & اِحْمَالَّ & حَمَرَ & اِحْمَارَّ\\
\vgrey & 19 & فَعَلَ & اِفْعَوَّلَ & حَمَلَ & اِحْمَوَّلَ & جَلَذَ & اِجْلَوَّذَ\\
\vlabel{6}{ثُلَاثِىْ مَزِيْدٌ}{بِثَلَاثَةِ أَحْرُفٍ} & 20 & فَعَلَ & اِفْعَوْعَلَ & حَمَلَ & اِحْمَوْمَلَ & خَشَنَ & اِخْشَوْشَنَ\\\hline
\hlabel{رُبَاعِىْ}{مُجَرَّدٌ} & 21 & فَعْلَلَ & فَعْلَلَ & دَحْرَجَ & دَحْرَجَ & زَلْزَلَ & زَلْزَلَ\\\hline
\hlabel{رُبَاعِىْ مَزِيْدٌ}{بِحَرْفٍ} & 22 & فَعْلَلَ & تَفَعْلَلَ & دَحْرَجَ & تَدَحْرَجَ & سَرْوَلَ & تَسَرْوَلَ\\\hline
\vgrey & 23 & فَعْلَلَ & اِفْعَنْلَلَ & دَحْرَجَ & اِدْحَنْرَجَ & حَرْجَمَ & اِحْرَنْجَمَ\\
\vlabel{2}{رُبَاعِىْ مَزِيْدٌ}{بِحَرْفَيْنِ} & 24 & فَعْلَلَ & اِفْعَلَلَّ & دَحْرَجَ & اِدْحَرَجَّ & قَشْعَرَ & اِقْشَعَرَّ\\
\end{sarf}
